\section{Introduction}
\label{sec:introduction}

Personalization is moving from a transient prompt into persistent model state. A mobile assistant may retrieve an old message, cache a preference as memory, or update a compact adapter with examples selected from recent interactions. Parameter-efficient adaptation makes the last option technically plausible: LoRA confines updates to low-rank matrices \cite{hu2021lora}, QLoRA lowers memory demand further \cite{dettmers2023qlora}, and prototype systems now demonstrate local fine-tuning on phones \cite{peng2024pocketllm,geng2025mobilefinetuner}. Personalization benchmarks such as LaMP formalize the utility side of this trend \cite{salemi2023lamp}. These capabilities are useful precisely because a phone sees data that a server does not.

The same locality creates an unusually sharp integrity problem. A phone sees user-authored text, but it also sees remote content delivered as messages, notifications, invitations, shared files, clipboard transfers, and OCR-indexed screenshots. The operating system can often distinguish those origins at ingestion time. Once content is converted to a plain training example or summarized into a memory record, however, the distinction can disappear. A received message can consequently influence a reply model as if it were a user preference. If that memory later becomes an adapter-training example, deleting the message or memory no longer removes the learned behavior.

This paper calls that transition \emph{cross-sink laundering}. The term is narrower than generic data poisoning. It denotes a causal path from an untrusted source, through one persistent learning sink, into another sink whose state survives cleanup of the first. AgentPoison shows that poisoned memories or knowledge bases can steer retrieval-augmented agents \cite{chen2024agentpoison}; MINJA demonstrates query-only memory injection \cite{dong2025minja}; recent work studies stealthy email injection, self-reinforcing memories, and poisoned experience retrieval \cite{zhang2026memghost,yang2026zombie,srivastava2025memorygraft}. SkillJack is particularly close conceptually: an agent can transform poisoned trajectories into reusable skills that outlive source deletion \cite{ying2026skilljack}. PoisonGuard asks the mobile-learning analogue: what happens when a memory record becomes a gradient update to a locally trained LoRA adapter?

Existing defenses do not directly enforce this boundary. Memory-only filters can reject or down-rank suspicious records \cite{sunil2026memorypoison,wei2025amemguard}, and isolation can stop external data from entering a privileged agent context \cite{wen2026agentsys}. Textual backdoor detectors inspect inputs or behavior \cite{qi2020onion,yang2021rap}, while spectral or activation statistics inspect learned representations \cite{tran2018spectral,chen2018activation}. Federated robust aggregation compares updates from multiple clients \cite{blanchard2017krum,yin2018byzantine,pillutla2019rfa,cao2020fltrust}, but a single phone has no worker population or trusted server reference. Approximate unlearning is also an unsafe synonym for restoration: empirical evidence shows that practical unlearning methods may leave poisoning effects behind \cite{pawelczyk2025unlearning}. These methods remain valuable baselines, but none answers which downstream state consumed a later-revoked local item.

\PG treats provenance as a control-plane object rather than a descriptive log. Every ingested event carries an envelope containing channel class, application identity, remote principal when available, user-authorship evidence, capture time, integrity metadata, and derivation parents. A policy returns one of three decisions for each sink: admit, quarantine, or deny. Memory and adapter training consult the same provenance graph. Candidate adapters are never promoted directly; they are trained against a versioned window, subjected to behavioral and update-level audits, and promoted only when the policy and audit gate both pass. Each promoted checkpoint records the exact source-event closure that influenced it.

Revocation therefore becomes graph invalidation followed by rollback-and-replay. When an item is found malicious, PoisonGuard identifies all descendants, chooses the newest checkpoint outside the invalidated closure, restores it, and replays later admissible events with pinned code, model, optimizer, tokenizer, and randomness. This is more expensive than deleting a row but makes the recovery claim testable. We state the conditions under which the recovered memory and adapter equal the counterfactual state that would have arisen had the poison never been admitted. We also identify the conditions under which this statement fails.

\paragraph{Contributions.} This manuscript develops five contributions locked to the original PoisonGuard concept.

\begin{enumerate}
  \item We formalize a single-device, multi-channel threat model for mixed memory-plus-adapter personalization and define cross-sink laundering, including persistence after source deletion.
  \item We design a provenance-carried selection plane that preserves source and derivation information from mobile ingestion through memory construction, training-window selection, candidate adapters, and promoted checkpoints.
  \item We introduce quarantine-then-promote adapter training and a shared lineage store that supports causal invalidation across both persistent sinks.
  \item We give sufficient-condition theorems for non-laundering and exact rollback-and-replay, together with explicit bounds on contaminated influence, checkpoint storage, and replay work. The claims are conditional, not optimality results.
  \item We specify an implementable evaluation on Android-class phones, including channel-adapted memory and LoRA poisoning, an adaptive user-retyping attacker, and systems metrics. No fabricated device results are reported.
\end{enumerate}

\paragraph{Scope and non-goals.} PoisonGuard does not protect a compromised OS or training runtime. It does not prove that content is semantically truthful. It does not solve generic prompt injection, privacy leakage, or federated Byzantine aggregation. It cannot distinguish a genuine user choice from a user action induced by social engineering without additional confirmation. Its strongest result is counterfactual restoration under explicit reproducibility assumptions, not universal prevention.

\paragraph{Road map.} \Cref{sec:related} positions the work. \Cref{sec:preliminaries,sec:system} define notation and the system. \Cref{sec:problem} formalizes laundering and objectives. \Cref{sec:method} presents the design, \cref{sec:theory} analyzes it, and \cref{sec:algorithms} gives executable procedures. \Cref{sec:experiments} specifies the evaluation; \cref{sec:discussion,sec:limitations} discuss deployment and limitations.

\begin{figure*}[t]
  \centering
  \resizebox{0.95\textwidth}{!}{\resizebox{\linewidth}{!}{%
\begin{tikzpicture}[node distance=7mm and 8mm, every node/.style={transform shape}]
  \node[honest, text width=18mm] (typed) {\tikz\draw[pgblue,line width=.7pt] (0,0)--(.18,.18)--(.55,-.19);\\Typed input};
  \node[attack, text width=22mm, below=5mm of typed] (remote) {\tikz{\draw[pgorange] (0,0) rectangle (.5,.32); \draw[pgorange] (0,.32)--(.25,.12)--(.5,.32);}\\Messages, apps, files};
  \node[defense, text width=23mm, right=12mm of typed, yshift=-7mm] (tag) {Provenance\\tagger};
  \node[defense, text width=23mm, right=of tag] (policy) {Sink-specific\\policy};
  \node[honest, text width=21mm, above right=7mm and 10mm of policy] (memory) {\tikz{\draw[pgblue] (-.25,-.2) rectangle (.25,.2); \draw[pgblue] (-.25,.05)--(.25,.05) (-.25,-.08)--(.25,-.08);}\\Memory store};
  \node[infra, text width=21mm, below right=7mm and 10mm of policy] (trainer) {Candidate LoRA\\trainer};
  \node[defense, text width=20mm, right=10mm of trainer] (audit) {Audit and\\promotion};
  \node[honest, text width=20mm, above=11mm of audit] (active) {Active\\adapter};
  \node[defense, text width=22mm, below=11mm of audit] (lineage) {\tikz\draw[pggreen] (0,0) circle (.09) (.38,0) circle (.09) (.19,-.27) circle (.09) (0,0)--(.38,0)--(.19,-.27)--cycle;\\Lineage store};

  \draw[dataflow] (typed.east) -| (tag.north west);
  \draw[attackflow] (remote.east) -| (tag.south west);
  \draw[dataflow] (tag.east) -- (policy.west);
  \draw[dataflow] (policy.east) -| (memory.west);
  \draw[dataflow] (policy.east) -| (trainer.west);
  \draw[dataflow] (memory.south) |- (trainer.north);
  \draw[dataflow] (trainer.east) -- (audit.west);
  \draw[controlflow] (audit.north) -- node[flowlabel]{promote} (active.south);
  \draw[controlflow] (tag.south) |- (lineage.west);
  \draw[controlflow] (policy.south) |- (lineage.west);
  \draw[controlflow] (memory.east) -| (lineage.north);
  \draw[controlflow] (audit.south) -- (lineage.north);

  \node[group, fit=(typed)(remote), label={[font=\scriptsize,text=pggray]above:Mobile ingestion}] {};
  \node[group, fit=(tag)(policy)(memory)(trainer)(audit)(active)(lineage), label={[font=\scriptsize,text=pggray]above:Trusted personalization boundary}] {};
  \coordinate (leg) at ($(lineage.south)+(0,-8mm)$);
  \draw[dataflow] ($(leg)+(-38mm,0)$)--($(leg)+(-30mm,0)$);
  \node[font=\scriptsize, anchor=west] at ($(leg)+(-28mm,0)$) {data/model};
  \draw[attackflow] ($(leg)+(-5mm,0)$)--($(leg)+(3mm,0)$);
  \node[font=\scriptsize, anchor=west] at ($(leg)+(5mm,0)$) {untrusted};
  \draw[controlflow] ($(leg)+(26mm,0)$)--($(leg)+(34mm,0)$);
  \node[font=\scriptsize, anchor=west] at ($(leg)+(36mm,0)$) {policy/recovery};
\end{tikzpicture}%
}
}
  \caption{PoisonGuard architecture. Blue solid arrows carry ordinary data, orange dashed arrows show potential adversarial flow, and green double-line arrows denote policy or recovery control. Provenance remains attached across both memory and adapter sinks.}
  \label{fig:architecture}
\end{figure*}
